\documentclass[10pt,a4paper]{amsart}
\usepackage[margin=19mm]{geometry}
\usepackage{amsmath,amssymb,mathtools}
\usepackage[scr=rsfs,cal=euler]{mathalfa}
\usepackage{xfrac}
\usepackage[unicode,hidelinks,bookmarks=false]{hyperref}
\newcommand{\ie}{i.e.\ }
\newcommand{\cf}{cf.\ }
\newcommand{\resp}{respectively\ }
\newcommand{\Subsection}{\subsection}
\providecommand{\nobreakdash}{\nobreak\hbox{-}}
\setlength{\emergencystretch}{2em}
\multlinegap0pt
\usepackage{bm}
\usepackage{bbm}
\usepackage{dsfont}
\usepackage{xspace}
\usepackage{cancel}
\usepackage{mathtools}
\usepackage{amsthm}
\makeatletter
\@ifundefined{thm}{\newtheorem{thm}{Thm}}{}
\@ifundefined{lemma}{\newtheorem{lemma}[thm]{Lemma}}{}
\makeatother
\newcommand{\Kk}{\mathcal {K}}
\renewcommand{\div}{\operatorname{div}}
\title[Graphical Willmore Problems with Low-Regularity Boundary and]{Graphical Willmore Problems with Low-Regularity Boundary and Dirichlet Data}
\author{Boris Gulyak}
\date{Source: arXiv:2509.21018v1 (CC BY 4.0)}
\begin{document}
\maketitle

\noindent\emph{Source and licence.} Graphical Willmore Problems with Low-Regularity Boundary and Dirichlet Data. Version arXiv:2509.21018v1 (CC BY 4.0), distributed under the Creative Commons Attribution 4.0 International licence, as recorded in the arXiv licence metadata for this exact version and shown on the abstract page https://arxiv.org/abs/2509.21018v1. Full rights notice: licenses/arxiv-2509.21018-CC-BY-4.0.txt.\par\smallskip

\noindent\textit{Editorial scope.} Counted unit: the Lemma whose source label is \texttt{lem:willrewriten} in the article (source0.tex lines 268--283, source label \texttt{lem:willrewriten}), delivered together with the article's own complete original proof of it (source0.tex lines 284--351, one \texttt{proof} environment of 68 lines). No mathematics is written, repaired or re-derived: statement and proof are the article's own text line for line. Required context retained uncounted: eq:willmoreeq (lines 247--250); eq:Hstar (lines 258--260). Every definition, display and label of those lines is retained unchanged. Omitted from this package and not redistributed: 1--246, 251--257, 261--267, 352--1344. Nothing inside a retained excerpt is omitted or altered: every retained body is delivered exactly as the article's lines are written, so no omission receipt of any kind accompanies this package. Citations: the retained excerpts carry no citation command at all, so no bibliography entry of the article is transcribed and no reference is redistributed. Reference adaptation: none was needed and none was made; every reference inside the delivered unit resolves inside it, so no unresolved reference or citation is produced. Printed slips kept literal: no printed word, symbol, hypothesis or number of the counted statement or of its proof was altered. Layout and class: the article's own class is \texttt{article} with packages including inputenc, amssymb, amsthm, babel, amsmath, mathtools, bbm, geometry, float, fancyhdr, hyperref, bookmark, tikz, enumerate; the wrapper is the lane's pinned \texttt{amsart} editorial wrapper at 10pt on A4, which restates the theorem-like environments the retained text needs and adds the article's own macro definitions that the retained text needs (\texttt{\textbackslash Kk}, \texttt{\textbackslash div}), with extra wrapper packages (bm, bbm, dsfont, xspace, cancel, mathtools). The delivered numbering is the unit's own. Assets: no retained span contains a figure environment, a graphics-inclusion command, a PDF-page inclusion, a table or a TikZ picture; no image file is copied into this package, the package declares no asset dependency and no image file is redistributed with it. Licence: version arXiv:2509.21018v1 of this article is distributed under the Creative Commons Attribution 4.0 International licence, as recorded in the arXiv licence metadata for this exact version and shown on the abstract page https://arxiv.org/abs/2509.21018v1. A full rights notice is delivered at licenses/arxiv-2509.21018-CC-BY-4.0.txt. Characters: every retained excerpt is pure ASCII, so the delivered engine, XeLaTeX with its default Latin Modern text font, reports no missing character.
\begin{align}
  0 =  \Delta_g H + \frac{1}{2}H^3-2H \Kk= \div \left( \frac{1}{Q} \left( \left(I-\frac{\nabla u \otimes \nabla u }{Q^2}
\right) \nabla (Q H) \right) - \frac{H^2 }{2Q} \nabla u\right), \label{eq:willmoreeq}
\end{align}

\begin{align}
    H=   Q^{-1}\star D^2u + Q^{-3}\star\nabla u\star D^2 u \star \nabla u. \label{eq:Hstar}
\end{align}

 \begin{lemma} \label{lem:willrewriten}
 The graphical Willmore equation \eqref{eq:willmoreeq} can be reformulated in the following divergence form
 \begin{align}
     0= \div \left( \frac{1}{Q} \left( \left(I-\frac{\nabla u \otimes \nabla u }{Q^2}
\right) \nabla (Q H) \right) - \frac{H^2 }{2Q} \nabla u\right)
= \Delta^2 u - D_{i} b^{i}_1[u] - D^2_{ij} b^{ij}_2[u],  \label{eq:Willmoreellipt}
\end{align}
where the lower-order terms are given by
\begin{align}
\begin{aligned}
b_1[u] =&\ D^2 u \star D^2 u \star \sum_{k=1}^3 Q^{-2k-1}P_{2k-1}(\nabla u),   \\
    b_2[u]=&\ D^2 u \star \sum_{k=1}^2 Q^{-2k-1}P_{2k}(\nabla u) + D^2 u \star P_2 (\nabla u) \star (Q(1+Q))^{-1}. 
\end{aligned}
    \label{eq:bi}
\end{align}
 \end{lemma}

\begin{proof}
We reformulate the Willmore equation by analyzing its terms individually. We begin by computing the gradient of the scalar function $Q$, which appears in the first term containing the biharmonic component:
\begin{align*}
    \nabla (Q)= \frac{D^2u \, \nabla u}{Q}.
\end{align*}
This identity is needed to handle the divergence term involving $QH$
\begin{align*}
    \div\left( \frac{1}{Q} \nabla (Q H) \right) 
    =&\  \div\left(  \nabla ( H) + \frac{D^2u \nabla u}{Q^3} QH \right) \\
    = &\ \Delta \left( \frac{\Delta u }{Q} \right) - \Delta \left( \frac{\nabla u (D^2 u \nabla u)}{Q^3} \right)  +\div\left( \frac{D^2u \nabla u}{Q^3} QH  \right).
\end{align*}
Especially since $(1-Q)(1+Q)=1-Q^2 = 1-1-|\nabla u|^2= -|\nabla u|^2$, it follows
\begin{align*}
\Delta \left( \frac{\Delta u }{Q} \right)   =& \Delta \left( \Delta u+\left(\frac{1}{Q}-1\right)\Delta u \right) = 
\Delta ^2 u +\Delta \left(\frac{1-Q}{Q}\Delta u \right) =
\Delta ^2 u - \Delta \left( \frac{|\nabla u|^2}{Q(1+Q)}\Delta u  \right).
\end{align*}
Combining all terms, we arrive at the identity
\begin{align*}
        \div\left( \frac{1}{Q} \nabla (Q H) \right) 
    = \Delta ^2 u - \Delta \left( \frac{|\nabla u|^2}{Q(1+Q)}\Delta u  \right) - \Delta \left( \frac{\nabla u (D^2 u \nabla u)}{Q^3} \right)  +\div\left( \frac{D^2u \nabla u}{Q^3} QH  \right).
\end{align*}
The second term in \eqref{eq:willmoreeq} can be rewritten using Einstein summation notation. Specifically, we consider
\begin{align*}
 \nabla_i\Bigg( & \frac{\nabla_i u \nabla _j u}{Q^3}  \nabla_j (Q H) \Bigg) 
    = \nabla_i\left( \nabla_j \left( \frac{\nabla_i u \nabla _j u}{Q^3}Q H \right)
     -\nabla_j\left( \frac{\nabla_i u \nabla _j u}{Q^3} \right)  Q H 
     \right)  \\
     =&\  \nabla_i\left( \nabla_j \left( \frac{\nabla_i u \nabla _j u}{Q^3}Q H \right)
     -(-3)\nabla_j (Q) \frac{\nabla_i u \nabla _j u}{Q^4}   Q H 
     - \frac{\nabla_j(\nabla_i u \nabla _j u)}{Q^3}   Q H 
     \right)  \\
     =&\  \nabla_i\left( \nabla_j \left( \frac{\nabla_i u \nabla _j u}{Q^3}Q H \right)
     + 3  \frac{\nabla_i u \nabla _j u (D^2u \nabla u)_j }{Q^5}Q H  -   \frac{D^2_{ij} u \nabla _j u +D^2_{jj} u \nabla _i u }{Q^3}Q H 
     \right)  \\
    =&\ D_{ij}^2 \left( \frac{\nabla_i u \nabla _j u}{Q^2} H \right)
     +  \div  \left( 3 \frac{\nabla u\cdot ( D^2 u \nabla u) }{Q^4}H \nabla u - \frac{H}{Q^2}  D^2u \nabla u  -   \frac{ \Delta u H  }{Q^2}   \nabla  u
     \right) \\
     =&\ D_{ij}^2 \left( \frac{\nabla_i u \nabla _j u}{Q^2} H \right)
     +  \div  \left( -3\frac{H^2 }{Q} \nabla u - \frac{H}{Q^2}  D^2u \nabla u  +2  \frac{ \Delta u H  }{Q^2}   \nabla  u
     \right).
\end{align*}
Combining the computations above, we arrive at the following representation:
\begin{align*}
    \Delta_g H + \frac{1}{2}H^3-2H \Kk
    = &\ \Delta ^2 u - \Delta \left( \frac{|\nabla u|^2}{Q(1+Q)}\Delta u +\frac{\nabla u (D^2 u \nabla u)}{Q^3} \right)   -D_{ij}^2 \left( \frac{\nabla_i u \nabla _j u}{Q^2} H \right) \\
 &\ 
     +  \div  \left( \frac{5}{2}\frac{H^2 }{Q} \nabla u + 2\frac{H}{Q^2}  D^2u \nabla u  -2  \frac{ \Delta u H  }{Q^2}   \nabla  u
      \right).
\end{align*}
Using the star notation introduced in \eqref{eq:Hstar}, where $ H=  D^2 u \star \sum_{k=1}^2 Q^{-2k+1}P_{2k-2}(\nabla u)$, we can express the right-hand side in an abstract polynomial form. This yields the final reformulation of the Willmore equation:
\begin{align*}
    \Delta_g H + \frac{1}{2}H^3-2H \Kk 
    = &\ \Delta ^2 u + \Delta \Big(D^2u \star P_2 (\nabla u) \star \big(Q(1+Q)\big)^{-1} \Big) + \Delta \big( D^2u\star Q^{-3}P_2(\nabla u) \big)\\
   &\  + D^2 \left(  D^2 u \star \sum_{k=1}^2 Q^{-2k-1}P_{2k}(\nabla u)\right) \\
    &\
     +  \div  \Big(  D^2 u \star D^2 u \star \nabla u \star Q^{-1}\big(
     Q^{-2} + Q^{-4}P_{2}(\nabla u) +Q^{-6}P_{4}(\nabla u) \big) 
      \Big)\\
    &\ +\div  \left(  D^2 u \star D^2 u \star \sum_{k=1}^2 Q^{-2k-1}P_{2k-1}(\nabla u) 
      \right).
\end{align*}
Altogether, the graphical Willmore equation can now be written as
\begin{align*}
     0 = \Delta_g H + \frac{1}{2}H^3-2H \Kk= \Delta^2 u - D_{i} b^{i}_1[u] - D^2_{ij} b^{ij}_2[u],  
\end{align*}
where the divergence terms $b^{i}_1[u]$ and $b^{ij}_2[u]$ are composed of combinations of lower-order terms as expressed above.
\end{proof}

\end{document}
